% !TEX root = ../main.tex
\questiontitle{Question 2: IEEE 802.11 RTS/CTS (Request to Send / Clear to Send) Mechanism}

\begin{question}{2}
Consider a smart agriculture IoT system in which battery-powered soil-moisture sensors A and B
periodically send their sensed data to a Wi-Fi (Wireless Fidelity)-enabled IoT Access Point (AP).
Due to the field layout and obstacles, A and B cannot hear each other directly, although both are within
communication range of the AP (Access Point).
\end{question}

% ------------------------------------------------------------------ (a)
\begin{question}{2 (a)}
Explain why a collision may occur when A and B attempt to send their sensor readings to the AP (Access Point) at
approximately the same time.
\end{question}

\answerhead{(a) Why can a collision occur?}
Sensors A and B cannot hear each other, but both can communicate with the AP (Access Point). If both transmit
simultaneously, their signals may overlap at the AP (Access Point), causing a packet collision. This is an example of
the hidden-terminal problem: A and B cannot detect each other's transmissions, even though both share the same
receiver.

\simvideo{q2_hidden}{5}{0.9\linewidth}{hidden-terminal collision without RTS/CTS (Request to Send / Clear to Send)}

\newpage
% ------------------------------------------------------------------ (b)
\begin{question}{2 (b)}
Using the RTS/CTS (Request to Send / Clear to Send) mechanism of IEEE 802.11, describe the complete communication sequence when sensor A wants to transmit its soil-moisture
reading to the AP (Access Point).
\end{question}

\answerhead{(b) Complete RTS/CTS (Request to Send / Clear to Send) communication sequence}
When sensor A wants to send a soil-moisture reading, the sequence is:

\begin{enumerate}[label=\textbf{\arabic*.},leftmargin=1.6em,itemsep=5pt]
\item \textbf{Channel check}\\
Sensor A senses the wireless channel and waits until it is free.
\item \textbf{RTS (Request to Send)}\\
A sends an RTS (Request to Send) frame to the AP (Access Point), requesting permission to transmit and indicating
the expected reservation duration.
\item \textbf{CTS (Clear to Send)}\\
After receiving the RTS (Request to Send), the AP (Access Point) responds with a CTS (Clear to Send) frame, granting
permission and announcing the remaining reserved transmission duration.
\item \textbf{DATA transmission}\\
A sends its soil-moisture data to the AP (Access Point).
\item \textbf{ACK (Acknowledgement)}\\
The AP (Access Point) sends an ACK (Acknowledgement) after successfully receiving the data frame.
\end{enumerate}

The basic exchange is:

\begin{center}
\begin{tikzpicture}[
  box/.style={draw=black,line width=0.8pt,minimum width=2.75cm,minimum height=1.25cm,align=center,
              font=\bfseries\large},
  arr/.style={-{Stealth[length=3mm]},line width=1pt}]
\node[box] (rts) {RTS\\[-1pt]{\scriptsize\mdseries (Request to Send)}};
\node[box,right=0.9cm of rts] (cts) {CTS\\[-1pt]{\scriptsize\mdseries (Clear to Send)}};
\node[box,right=0.9cm of cts] (data) {DATA\\[-1pt]{\scriptsize\mdseries (soil reading)}};
\node[box,right=0.9cm of data] (ack) {ACK\\[-1pt]{\scriptsize\mdseries (Acknowledgement)}};
\draw[arr] (rts) -- (cts);
\draw[arr] (cts) -- (data);
\draw[arr] (data) -- (ack);
\tikzset{flowcap/.style={font=\scriptsize,text width=2.9cm,align=center}}
\node[flowcap,below=2pt of rts] {sensor A $\rightarrow$\\AP (Access Point)};
\node[flowcap,below=2pt of cts] {AP (Access Point) $\rightarrow$\\sensor A\\(B also hears it)};
\node[flowcap,below=2pt of data] {sensor A $\rightarrow$\\AP (Access Point)};
\node[flowcap,below=2pt of ack] {AP (Access Point) $\rightarrow$\\sensor A};
\end{tikzpicture}
\end{center}

Short interframe spaces (SIFS) separate the CTS (Clear to Send), DATA, and ACK (Acknowledgement) exchanges as
specified by IEEE 802.11.

\simvideo{q2_rtscts}{5}{0.9\linewidth}{complete RTS (Request to Send) / CTS (Clear to Send) / DATA / ACK (Acknowledgement) exchange}


\newpage
% ------------------------------------------------------------------ (c)
\begin{question}{2 (c)}
Explain how the CTS (Clear to Send) message received by sensor B helps prevent B from transmitting while A is
communicating with the AP (Access Point).
\end{question}

\answerhead{(c) How does CTS (Clear to Send) prevent sensor B from transmitting?}
Sensor B receives the AP's (Access Point's) CTS (Clear to Send) frame even though it cannot hear sensor A. The CTS
(Clear to Send) contains a duration field that informs B how long the medium is reserved. B updates its Network
Allocation Vector (NAV) and defers transmission during the reserved period.

Consequently, B avoids transmitting while A communicates with the AP (Access Point), reducing the likelihood of a
collision.

\simfigure[width=0.86\linewidth,page=41]{../media/frames/q2_rtscts.pdf}{frame from the video above: B's NAV (Network Allocation Vector) keeps B silent while A sends DATA}

\newpage
% ------------------------------------------------------------------ (d)
\begin{question}{2 (d)}
From an IoT perspective, explain one benefit and one drawback of using RTS/CTS (Request to Send
/ Clear to Send) for battery-powered sensor nodes.
\end{question}

\answerhead{(d) One benefit and one drawback for battery-powered IoT nodes}
\textbf{Benefit:} RTS/CTS (Request to Send / Clear to Send) reduces hidden-terminal collisions, avoiding wasted
energy from unsuccessful data transmissions and retransmissions.

\textbf{Drawback:} RTS (Request to Send) and CTS (Clear to Send) add control-frame overhead, extra channel access
time, and additional energy consumption. For small, frequent sensor packets, this overhead may outweigh the
benefits.

